% Brief (grammar: a small table of the possible combinations, with a mark
%   on each one that occurred; atlas VI/fig:bonded-pairs).
% Reader question: which combinations of a claim's status and its bond's
% status occur, in a harbor where a claim requires a bond?
% Claim: in the kernel's harbor four combinations occur, and in each the
% bond is held while the claim is not released. A harbor that decides
% claims and funds separately reaches six more with its first wrong
% operation, each a claim with no held bond, and it admits writes after a
% bond is refunded.
% Evidence: every state each harbor reached in the exhaustive test
% [internal, harbor stakes --each; sessions stakes-each.txt and
% stakes-each-separate.txt, drawn as printed, and
% tests/harbor-research/test_bonded_stake_numbers.py compares the counts
% in the caption with them].
% Must distinguish: a combination the reference model allows from one it
% does not; a claim under which writes were admitted from one under which
% none were.
% Grammar chosen: the claim's status in rows and the bond's in columns, so
% that the allowed combinations are a fixed set of cells, which are tinted,
% and a mark outside the tint is visibly a fault.
% Grammar rejected: a state diagram of one claim, which shows the order of
% operations and hides which combinations occur; bars of the counts, whose
% sizes depend on the bounds and mean nothing.
% Five-second test: in which table is there a mark outside the tint?
\begin{figure}[H]
\centering
\pdfigurehue{pdcobalt}%
\begin{tikzpicture}[pd figure,x=1.4cm,y=.6cm,
    allowed/.style={fill=pdfocus!11},
    occurs/.style={circle,fill=pdink,inner sep=0pt,minimum size=4.4pt},
    unbonded/.style={circle,fill=pderror,inner sep=0pt,minimum size=4.4pt},
    written/.style={draw=pdink,line width=.5pt},
    written unbonded/.style={draw=pderror,line width=.8pt}]
  % The tables are the recorded runs, read as printed: the first line is
  % the command, and every comment line starts with '#'.
  \pgfplotstableread[skip first n=1]{excerpts/stakes-each.txt}\pdpairsreal
  \pgfplotstableread[skip first n=1]{excerpts/stakes-each-separate.txt}\pdpairsseparate
  % Counts computed from the tables as they are drawn.
  \ifdefined\pdpairsin\else \newcount\pdpairsin \newcount\pdpairsout \fi
  % One panel. A claim's status is a row, with live at the top. A bond's
  % status is a column.
  \newcommand{\pdpairspanel}[1]{%
    \global\pdpairsin=0 \global\pdpairsout=0
    % The combinations the model allows: a held bond with a claim that is
    % not released, and a refunded bond with one that is.
    \foreach \y in {3,2,0}{ \fill[allowed] (.5,\y-.5) rectangle (1.5,\y+.5); }
    \fill[allowed] (1.5,.5) rectangle (2.5,1.5);
    \foreach \x in {-.5,.5,1.5,2.5,3.5}{ \draw[pdinkmuted!35,line width=.3pt] (\x,-.5) -- (\x,3.5); }
    \foreach \y in {-.5,.5,1.5,2.5,3.5}{ \draw[pdinkmuted!35,line width=.3pt] (-.5,\y) -- (3.5,\y); }
    \pgfplotstablegetrowsof{#1}
    \pgfmathtruncatemacro{\pdlastpair}{\pgfplotsretval-1}
    \foreach \i in {0,...,\pdlastpair}{
      \pgfplotstablegetelem{\i}{claim}\of#1 \let\cl\pgfplotsretval
      \pgfplotstablegetelem{\i}{bond}\of#1 \let\bo\pgfplotsretval
      \pgfplotstablegetelem{\i}{writes}\of#1 \let\ch\pgfplotsretval
      \pgfmathtruncatemacro{\y}{3-\cl}
      % A claim that is not released, with a bond that is not held.
      \pgfmathtruncatemacro{\nobond}{and(\cl!=2,\bo!=1)}
      \ifnum\nobond=1 \global\advance\pdpairsout 1 \node[unbonded] at (\bo,\y) {};
      \else \global\advance\pdpairsin 1 \node[occurs] at (\bo,\y) {};\fi
      % A ring: writes were admitted under such a claim.
      \ifnum\ch>0
        \ifnum\bo=1 \draw[written] (\bo,\y) circle[radius=4.6pt];
        \else \draw[written unbonded] (\bo,\y) circle[radius=4.6pt];\fi
      \fi
    }
    \foreach \name [count=\row from 0] in {live,expired,released,abandoned}{
      \node[pd label,text=pdinkmuted,anchor=east] at (-.6,3-\row) {\name};
    }
  }
  \pdpairspanel{\pdpairsreal}
  \foreach \name [count=\x from 0] in {none,held,refunded,slashed}{
    \node[pd label,text=pdinkmuted,anchor=base] at (\x,3.85) {\name};
  }
  \node[pd label,anchor=base east] at (-.6,3.85) {the claim};
  \node[pd label,anchor=base] at (1.5,4.85) {its bond};
  \node[pd title,anchor=north west,align=left] at (3.75,3.4) {the kernel's harbor};
  \node[pd label,anchor=north west,align=left] at (3.75,2.5)
    {\the\pdpairsin\ combinations, each\\in a tinted cell};
  \begin{scope}[yshift=-2.95cm]
    \pdpairspanel{\pdpairsseparate}
    \node[pd title,anchor=north west,align=left] at (3.75,3.4) {claims and funds\\decided separately};
    \node[pd label,anchor=north west,align=left] at (3.75,1.6)
      {\the\pdpairsout\ more, each a claim\\with no held bond};
  \end{scope}
\end{tikzpicture}
\caption{\textbf{In the kernel's harbor a claim and its bond occur in
four combinations.} A row is the claim's status and a column is the status
of its bond. A mark means that some state of the exhaustive test contained
that combination, and a ring means that writes were admitted under such a
claim. The tinted cells are the combinations the reference model allows.
In the lower table, claims and funds are decided by two components that
are each correct, and the red marks are claims with no held bond. The red
ring is writes admitted after the bond was refunded. [internal,
\texttt{harbor stakes -{}-each}]}
\label{fig:bonded-pairs}
\end{figure}
